(1) is the census (\claimref{sheaf.collision-census}) plus the admission of equality in Definition \claimref{obj.sheaf.variable-order-algebra-strict-diagonal-data}; the parametrisation identity and its fixed point are exact rational arithmetic. (2) is the census applied along the change, with every stated property of $q(v)$ an exact rational-function identity in $v$; note $q(\tfrac12)$ is the $u=2$ triple and $q(-\tfrac14)$ is the $u=3$ triple with axes $1,2$ exchanged, so the curve traverses genuine Kasner data on both sides. (3) The proof of Theorem \claimref{sheaf.variable-weight-boundary-tameness-completion-forced} used the ordering only to name the dying axis, so it applies on each chamber with the axes relabelled; the script verifies the full relabelled package on the $v<0$ chamber at $q(-\tfrac14)$: forced weights $(0,q_1,0,q_3)$, weighted Gram $\operatorname{diag}(1,-1,-\varrho',-1)$ with $\varrho' = \tau^{-2q_2}$, boundary corank one with radical on axis $2$, coupling determinant $1+\varrho'$, completion signature $(2,4)$ with no shadow wall, and the two weight-forcing solves. The dying-order identity is checked as a rational-function identity, and its Schwarzschild consistency as well: at $u=1$ the order $\tfrac23$ in $\tau$ converts to first order in $r$, by $\varrho^3 = \tau^2 = \tfrac49 r^3(1 + O(r))$ exactly, matching $\det G = -r$ of Theorem \claimref{sheaf.schwarzschild-degeneration-boundary-tame}. On $Z$ the weights $(0,0,0,1)$ rescale only the third axis and the four Gram entries are read off directly; the signature $(1,3,0)$ is checked at two unit choices, while the degeneration sheaf keeps its rank sixteen there by \claimref{sheaf.local-freeness-constant-rank}. (4) The coincidence and the two negative transition exponents are exact evaluations on the curve.
